\chapter{De ruimten \texorpdfstring{$L^p$}{Lp}}\label{ch:b3:lp}

De Lebesgue-integraal is voor de analyse gebouwd; de ruimten $L^p$
zijn de plek waar die analyse woont. Het zijn banachruimten
(Riesz--Fischer) --- de \omterm{thm:b3:complete:completion}{vervolledigingen} die de \omterm{def:b3:topology:continuity}{continue} functies
volgens \cref{exo:b3:complete:1} misten --- en ze dragen een
gladstrijktechniek, \emph{convolutie tegen \omterm{def:b3:lp:mollifier}{mollifiers}}, die elke
$L^p$-functie door $\mathcal C^\infty$-functies benadert. Dit
hoofdstuk bewijst de integraalversies van Hölder en Minkowski, de
\omterm{def:b3:complete:complete}{volledigheid}, de dichtheidsstellingen en de regularisatiemachine,
en eindigt met de inclusie- en interpolatiegeografie van de
$L^p$-schaal. Overal is $(X, \mathcal A, \mu)$ een maatruimte en
zijn de functies complexwaardig; op $\R^d$ is de \omterm{def:b3:measure:measure}{maat} $\lambda_d$.

\section{Definitie; Hölder en Minkowski}

\begin{definition}\label{def:b3:lp:space}
Voor $1 \leq p < \infty$ is $\mathcal L^p(\mu)$ de verzameling
\omterm{def:b3:lebesgue:measurable}{meetbare} $f$ met $\norm f_p = \bigl(\int\abs
f^p\dd\mu\bigr)^{1/p} < \infty$, en $\mathcal L^\infty(\mu)$ de
verzameling functies die buiten een nulverzameling begrensd zijn,
met $\norm f_\infty$ het \emph{essentiële supremum} --- de
kleinste $M$ met $\abs f \leq M$ b.o. (het infimum wordt bereikt:
snijd de nulverzamelingen voor $M + \frac1n$). Omdat $\norm f_p =
0$ slechts $f = 0$ \emph{b.o.} afdwingt
(\cref{exo:b3:lebesgue:5}), stellen we
\[
L^p(\mu) = \mathcal L^p(\mu)/\{f = 0 \text{ b.o.}\} :
\]
de elementen zijn klassen van functies modulo nulverzamelingen,
en $\norm\cdot_p$ is een echte norm op
$L^p$.\index{Lp-ruimte@$L^p$-ruimte}
\end{definition}

\begin{theorem}[Ongelijkheid van Hölder]\label{thm:b3:lp:holder}
Zij $1 \leq p, q \leq \infty$ met $\frac1p + \frac1q = 1$
(toegevoegde exponenten). Voor \omterm{def:b3:lebesgue:measurable}{meetbare} $f, g$
geldt\index{ongelijkheid van Hölder}
\[
\norm{fg}_1 \leq \norm f_p\,\norm g_q ,
\]
met gelijkheid (voor $1 < p < \infty$, eindige normen, $f,g\ne0$)
dan en slechts dan als $\abs f^p$ en $\abs g^q$ b.o. evenredig
zijn.
\end{theorem}

\begin{proof}
De gevallen $\{p, q\} = \{1, \infty\}$ zijn rechtstreeks
($\abs{fg} \leq \norm g_\infty\abs f$ b.o.). Zij $1 < p <
\infty$; normeer $\norm f_p = \norm g_q = 1$ (homogeniteit; nul-
of oneindige normen zijn triviaal). De ongelijkheid van Young $ab
\leq \frac{a^p}p + \frac{b^q}q$ ($a, b \geq 0$; concaviteit van
$\ln$, zoals in \cref{pb:b3:banach:1}) geeft puntsgewijs $\abs{f
g} \leq \frac{\abs f^p}p + \frac{\abs g^q}q$; integreer:
$\norm{fg}_1 \leq \frac1p + \frac1q = 1$. Gelijkheid dwingt
b.o.\ gelijkheid in Young af, dat wil zeggen $\abs f^p = \abs
g^q$ b.o. (na de normering; terugschalen geeft evenredigheid).
\end{proof}

\begin{theorem}[Ongelijkheid van
Minkowski]\label{thm:b3:lp:minkowski}
Voor $1 \leq p \leq \infty$ is $\norm{f + g}_p \leq \norm f_p +
\norm g_p$.\index{ongelijkheid van Minkowski}
\end{theorem}

\begin{proof}
$p = 1, \infty$: de puntsgewijze respectievelijk b.o.\
driehoeksongelijkheid. Voor $1 < p < \infty$ nemen we aan dat
$\norm{f+g}_p < \infty$ (anders levert $\abs{f+g}^p \leq
2^{p-1}(\abs f^p + \abs g^p)$, uit de convexiteit van $t^p$, dat
het linkerlid eindig is zodra het rechterlid dat is). Dan is
\[
\norm{f{+}g}_p^p \leq \int\abs f\,\abs{f{+}g}^{p-1} +
\int\abs g\,\abs{f{+}g}^{p-1}
\leq \bigl(\norm f_p + \norm g_p\bigr)\,
\bigl\|\abs{f{+}g}^{p-1}\bigr\|_q
\]
volgens Hölder, en $\norm{\abs{f+g}^{p-1}}_q =
\norm{f+g}_p^{p/q}$ omdat $(p-1)q = p$; deel door
$\norm{f+g}_p^{p/q}$ (als dat niet nul is; anders triviaal) en
gebruik $p - \frac pq = 1$.
\end{proof}

\section{Volledigheid en haar metgezellen}

\begin{theorem}[Riesz--Fischer]\label{thm:b3:lp:complete}
Voor $1 \leq p \leq \infty$ is $L^p(\mu)$ een banachruimte.
Bovendien heeft elke in $L^p$ convergerende rij een deelrij die
\emph{\omterm{def:b3:lebesgue:l1}{bijna overal}} convergeert (met een $L^p$-dominant in het
geval $p < \infty$).\index{stelling van Riesz--Fischer}
\end{theorem}

\begin{proof}
$p = \infty$: een $\norm\cdot_\infty$-cauchyrij is, buiten één
enkele nulverzameling (de vereniging van aftelbaar veel),
uniform cauchy: ze convergeert daarbuiten uniform; klaar. Zij $p
< \infty$. Volgens \cref{exo:b3:complete:1}(b) volstaat het
absoluut convergente reeksen te sommeren: zij
$\sum\norm{f_k}_p = M < \infty$. Stel $G_n = \sum_{k \leq
n}\abs{f_k}$ en $G = \sum_k\abs{f_k}$ (puntsgewijs in
$[0,\infty]$): volgens Minkowski is $\norm{G_n}_p \leq M$, en de
monotone convergentie ($G_n^p \nearrow G^p$) geeft $\int G^p
\leq M^p$: dus $G < \infty$ b.o., zodat de reeks $\sum f_k(x)$
voor b.o.\ $x$ absoluut convergeert; noem de som $S(x)$
(willekeurige waarde op de nulverzameling). Dan is $\abs{S -
\sum_{k\leq n}f_k}^p \leq (2G)^p \in L^1$, en de gedomineerde
convergentie geeft $\norm{S - \sum_{k\leq n}f_k}_p \to 0$: de
reeks convergeert in $L^p$.

De uitspraak over de deelrij: als $f_n \to f$ in $L^p$, kies dan
$n_k$ met $\norm{f_{n_{k+1}} - f_{n_k}}_p \leq 2^{-k}$; de reeks
$\sum(f_{n_{k+1}} - f_{n_k})$ valt onder het vorige argument: ze
convergeert b.o.\ absoluut, gedomineerd door een $G \in L^p$,
dus $f_{n_k} \to f_{n_1} + \sum(\cdots)$ b.o., en die b.o.-limiet
moet (een vertegenwoordiger van) $f$ zijn (beide zijn
$L^p$-limieten). De dominant: $\abs{f_{n_k}} \leq \abs{f_{n_1}} +
G$.
\end{proof}

\begin{remark}\label{rem:b3:lp:modes}
$L^p$-convergentie impliceert geen b.o.-convergentie (de rij van
de \emph{schrijfmachine}, \cref{exo:b3:lp:3}), en evenmin
omgekeerd (wegglippende bulten): de twee convergentiewijzen zijn
alleen via deelrijen en dominantie verbonden. De tegenvoorbeelden
van \cref{exo:b3:lp:3} in het geheugen houden is het beste
vaccin.
\end{remark}

\section{Dichtheidsstellingen}

\begin{theorem}\label{thm:b3:lp:density}
Zij $1 \leq p < \infty$.
\begin{enumerate}
  \item De \omterm{def:b3:lebesgue:simple}{elementaire functies} (met dragers van eindige \omterm{def:b3:measure:measure}{maat})
        liggen \omterm{def:b3:topology:interior}{dicht} in $L^p(\mu)$.
  \item In $L^p(\R^d)$ liggen de \omterm{def:b3:topology:continuity}{continue} functies met \omterm{def:b3:topology:compact}{compacte}
        drager, $\mathcal C_c(\R^d)$, \omterm{def:b3:topology:interior}{dicht}.
  \item De translatie is \omterm{def:b3:topology:continuity}{continu} op $L^p(\R^d)$: schrijven we
        $\tau_hf = f(\cdot - h)$, dan is $\norm{\tau_hf - f}_p
        \to 0$ als $h \to 0$.
\end{enumerate}
Geen van de drie geldt voor $p = \infty$.
\end{theorem}

\begin{proof}
(1) Voor $f \geq 0$: de dyadische $s_n \nearrow f$ van
\cref{thm:b3:lebesgue:approximation} voldoen aan $\abs{f -
s_n}^p \leq f^p \in L^1$: gedomineerde convergentie. (Elke $s_n
\leq f$ ligt in $L^p$, en haar niveauverzamelingen hebben
eindige \omterm{def:b3:measure:measure}{maat} waar de waarde positief is: $\mu(s_n \geq c) \leq
c^{-p}\int f^p$.) Splits een algemene $f$ in vier
niet-negatieve delen.

(2) Volgens (1) volstaat het $\mathbf 1_A$ te benaderen, met $A$
borel en $\lambda_d(A) < \infty$. De regulariteit (het bewijs
verloopt als in \cref{thm:b3:measure:regularity}) geeft een
\omterm{def:b3:topology:compact}{compacte} $K \subseteq A \subseteq U$ met $U$ open en
$\lambda_d(U\setminus K) < \varepsilon$; de functie van Urysohn
\[
\varphi(x) = \frac{d(x, \R^d\setminus U)}{d(x, \R^d\setminus U)
+ d(x, K)}
\]
is \omterm{def:b3:topology:continuity}{continu}, gelijk aan $1$ op $K$ en aan $0$ buiten $U$, en kan
\omterm{def:b3:topology:compact}{compacte} drager krijgen (krimp $U$ eerst tot een begrensde \omterm{def:b3:topology:topology}{open
verzameling}). Dan is $\norm{\mathbf 1_A - \varphi}_p^p \leq
\lambda_d(U\setminus K) < \varepsilon$.

(3) Voor $g \in \mathcal C_c$: de uniforme \omterm{def:b3:topology:continuity}{continuïteit} geeft
$\norm{\tau_hg - g}_\infty \to 0$, met dragers in één vaste
\omterm{def:b3:topology:compact}{compacte} verzameling voor $\abs h \leq 1$: dus $\norm{\tau_hg -
g}_p \to 0$. Voor algemene $f$: kies $g \in \mathcal C_c$ met
$\norm{f - g}_p < \varepsilon$; dan is $\norm{\tau_hf - f}_p
\leq 2\norm{f - g}_p + \norm{\tau_hg - g}_p$
(translatie-invariantie van de norm).

Voor $p = \infty$: uniforme benadering van $\mathbf
1_{\intoo0\infty}$ door \omterm{def:b3:topology:continuity}{continue} functies is onmogelijk (de
sprong), en $\norm{\tau_h\mathbf 1_{\intoo0\infty} - \mathbf
1_{\intoo0\infty}}_\infty = 1$ voor $h \neq 0$.
\end{proof}

\section{Convolutie en regularisatie}

\begin{theorem}[Ongelijkheid van Young]\label{thm:b3:lp:young}
Zij $1 \leq p \leq \infty$, $f \in L^1(\R^d)$ en $g \in
L^p(\R^d)$. Dan is $f * g$ b.o.\ gedefinieerd, behoort tot $L^p$
en geldt\index{ongelijkheid van Young}
\[
\norm{f * g}_p \leq \norm f_1\,\norm g_p .
\]
\end{theorem}

\begin{proof}
$p = \infty$: rechtstreekse afschatting. $p = 1$:
\cref{thm:b3:product:convolution}. Zij $1 < p < \infty$ en $q$
toegevoegd. Splits $\abs{f(y)} = \abs{f(y)}^{1/q}\cdot
\abs{f(y)}^{1/p}$ en pas Hölder toe:
\[
\int\abs{f(y)}\,\abs{g(x{-}y)}\,\dd y
\leq \Bigl(\int\abs f\Bigr)^{1/q}
\Bigl(\int\abs{f(y)}\,\abs{g(x - y)}^p\,\dd y\Bigr)^{1/p} .
\]
Verhef tot de $p$-de macht en integreer naar $x$; Tonelli op de
tweede factor geeft $\norm f_1^{p/q}\cdot\norm f_1\norm g_p^p$,
oftewel $\norm{f*g}_p^p \leq \norm f_1^{1 + p/q}\norm g_p^p =
(\norm f_1\norm g_p)^p$ --- en de eindigheid van de
Tonelli-integraal rechtvaardigt de b.o.\ absolute convergentie,
net als in \cref{thm:b3:product:convolution}.
\end{proof}

\begin{definition}[Mollifiers]\label{def:b3:lp:mollifier}
De functie\index{mollifier}\index{bultfunctie}
\[
\rho(x) = \begin{cases}
c\,\exp\Bigl(-\dfrac{1}{1 - \norm x^2}\Bigr) & \norm x < 1,\\
0 & \norm x \geq 1,
\end{cases}
\]
met $c$ zo genormeerd dat $\int\rho = 1$, is $\mathcal C^\infty$
op $\R^d$: de kern van de zaak is dat $t \mapsto
\eu^{-1/t}\mathbf 1_{t>0}$ $\mathcal C^\infty$ is op $\R$, met
alle afgeleiden in $0^+$ gelijk aan $0$ (elke afgeleide is
$P(1/t)\eu^{-1/t}$ voor een veelterm $P$, en dat gaat naar $0$;
inductie). Voor $\varepsilon > 0$ stellen we
$\rho_\varepsilon(x) = \varepsilon^{-d}\rho(x/\varepsilon)$: met
drager in $\bar B(0, \varepsilon)$ en nog steeds integraal $1$.
\end{definition}

\begin{theorem}[Regularisatie]\label{thm:b3:lp:regularization}
Zij $1 \leq p < \infty$ en $f \in L^p(\R^d)$. Dan geldt:
\begin{enumerate}
  \item $f * \rho_\varepsilon \in \mathcal C^\infty(\R^d)$, met
        $\partial^\alpha(f * \rho_\varepsilon) = f *
        \partial^\alpha\rho_\varepsilon$;
  \item $\norm{f * \rho_\varepsilon - f}_p \to 0$ als
        $\varepsilon \to 0$;
  \item bijgevolg ligt $\mathcal C^\infty_c(\R^d)$ \omterm{def:b3:topology:interior}{dicht} in
        $L^p(\R^d)$.
\end{enumerate}
\end{theorem}

\begin{proof}
(1) Differentiëren onder de integraal
(\cref{thm:b3:lebesgue:paramdiff}) naar $x$: voor $x$ in een bal
$B$ is $\abs{\partial_{x_i}\rho_\varepsilon(x - y)} \leq
C_\varepsilon\,\mathbf 1_{K}(y)$ met $K$ \omterm{def:b3:topology:compact}{compact} ($y$ op afstand
hoogstens $\varepsilon$ van $B$), en $\abs f\,\mathbf 1_K \in
L^1$ (Hölder tegen $\mathbf 1_K$): de stelling is van
toepassing; itereer voor hogere afgeleiden.

(2) Omdat $\int\rho_\varepsilon = 1$:
\[
(f * \rho_\varepsilon)(x) - f(x)
= \int \bigl(f(x - y) - f(x)\bigr)\rho_\varepsilon(y)\,\dd y ,
\]
en de \emph{integraalongelijkheid} van Minkowski --- of
rechtstreeks: Hölder/Jensen met de kansmaat
$\rho_\varepsilon\dd y$ en Tonelli --- geeft
\[
\norm{f*\rho_\varepsilon - f}_p^p
\leq \int\Bigl(\int\abs{f(x-y) -
f(x)}^p\dd x\Bigr)\rho_\varepsilon(y)\,\dd y
= \int \norm{\tau_yf - f}_p^p\;\rho_\varepsilon(y)\,\dd y
\]
(de tussenstap: pas de ongelijkheid van Jensen,
\cref{exo:b3:lp:10}, toe op de binnenste $y$-integraal en
vervolgens Tonelli). De integrand heeft drager in $\norm y \leq
\varepsilon$ en gaat daar uniform naar $0$ als $\varepsilon \to
0$ (\cref{thm:b3:lp:density}(3)): dus gaat de hele uitdrukking
naar $0$.

(3) Benader $f$ door $g \in \mathcal C_c$
(\cref{thm:b3:lp:density}(2)), en vervolgens $g$ door $g *
\rho_\varepsilon \in \mathcal C_c^\infty$ (\omterm{def:b3:topology:compact}{compacte} drager: de
som van de dragers).
\end{proof}

\begin{example}[$\abs x$ gladstrijken, met snelheden]
\label{ex:b3:lp:mollifyexample}
Neem $f(x) = \abs x$ op $\R$ (lokaal $L^1$; de stelling is op elk
begrensd venster van toepassing) en een symmetrische \omterm{def:b3:lp:mollifier}{mollifier}
$\rho_\varepsilon$. Dan is
\[
f_\varepsilon(x) = (f * \rho_\varepsilon)(x)
= \int\abs{x - y}\,\rho_\varepsilon(y)\,\dd y
\]
$\mathcal C^\infty$; ver van de knik gebeurt er niets: voor
$\abs x \geq \varepsilon$ is $\abs{x - y}$ lineair in $x$ op de
drager van $\rho_\varepsilon$, zodat $f_\varepsilon(x) = \abs x$
\emph{exact} (de symmetrie doodt de correctie). Nabij $0$ kost
het gladstrijken precies
\[
0 \leq f_\varepsilon(0) = \int\abs
y\,\rho_\varepsilon(y)\,\dd y \leq \varepsilon,
\qquad
\norm{f_\varepsilon - f}_\infty \leq \varepsilon :
\]
de benaderingsfout blijft opgesloten in de
$\varepsilon$-omgeving van de singulariteit en is van haar
grootte. Ondertussen is $f_\varepsilon'' \geq 0$ overal ($f$ is
convex, en convolutie tegen $\rho_\varepsilon \geq 0$ behoudt de
convexiteit), met $\int f_\varepsilon'' = f_\varepsilon'(\infty)
- f_\varepsilon'(-\infty) = 2$: de tweede afgeleide is een bult
van massa $2$ geperst in een breedte $O(\varepsilon)$, dus
$\norm{f_\varepsilon''}_\infty \gtrsim \varepsilon^{-1}$.
Gladstrijken is een ruil: uniforme fout $O(\varepsilon)$ tegen
explosie van de afgeleide $O(\varepsilon^{-1})$ --- precies de
wisselkoers die de kwantitatieve analyse
(interpolatieongelijkheden, de ideeënkring van
\cref{pb:b3:lp:1}) formaliseert.
\end{example}

\begin{corollary}[Hoofdlemma van de
variatierekening]\label{cor:b3:lp:fundlemma}
Zij $f \in L^1_{\mathrm{loc}}(\R^d)$ (\omterm{def:b3:lebesgue:l1}{integreerbaar} op \omterm{def:b3:topology:compact}{compacte}
verzamelingen) met $\int f\varphi = 0$ voor elke $\varphi \in
\mathcal C^\infty_c(\R^d)$. Dan is $f = 0$ b.o.\index{hoofdlemma
van de variatierekening}
\end{corollary}

\begin{proof}
Leg een bal $B = B(0, R)$ vast en stel $g = f\mathbf 1_{B(0,
R+1)} \in L^1$. Voor $x \in B$ en $\varepsilon < 1$ is $(g *
\rho_\varepsilon)(x) = \int f(y)\rho_\varepsilon(x - y)\dd y =
0$, want de testfunctie $y \mapsto \rho_\varepsilon(x-y)$ ligt in
$\mathcal C_c^\infty$. Maar $g * \rho_\varepsilon \to g$ in $L^1$
(\cref{thm:b3:lp:regularization}): dus $g = 0$ b.o.\ op $B$; put
$\R^d$ uit.
\end{proof}

\section{De geografie van \texorpdfstring{$L^p$}{Lp}}

\begin{proposition}\label{prop:b3:lp:inclusions}
(a) Is $\mu(X) < \infty$ en $1 \leq p \leq q \leq \infty$, dan is
$L^q \subseteq L^p$ met $\norm f_p \leq \mu(X)^{\frac1p -
\frac1q}\,\norm f_q$.
(b) Op $\R^d$ (oneindige \omterm{def:b3:measure:measure}{maat}) zijn er geen inclusies: voor $p
\neq q$ bestaan er functies in $L^p\setminus L^q$.
(c) (Interpolatie) Is $p < r < q$ en wordt $\alpha \in \intoo01$
bepaald door $\frac1r = \frac\alpha p + \frac{1 - \alpha}q$, dan
is
\[
\norm f_r \leq \norm f_p^{\alpha}\,\norm f_q^{1 - \alpha} ;
\]
in het bijzonder is $L^p \cap L^q \subseteq L^r$.
\end{proposition}

\begin{proof}
(a) Hölder met de exponenten $\frac qp$ en zijn toegevoegde:
$\int\abs f^p\cdot 1 \leq \norm{\abs f^p}_{q/p}\,\norm
1_{(q/p)'} = \norm f_q^p\,\mu(X)^{1 - p/q}$ (het geval $q =
\infty$ rechtstreeks). (b) Nabij $0$ en nabij $\infty$ ijken de
machten $x^{-\alpha}$ de zaak: \cref{exo:b3:lp:2}. (c) Schrijf
$\abs f^r = \abs f^{r\alpha}\abs f^{r(1-\alpha)}$ en pas Hölder
toe met het toegevoegde paar $\frac p{r\alpha}$, $\frac
q{r(1-\alpha)}$ (toegevoegd precies wegens de definitie van
$\alpha$): $\int\abs f^r \leq \norm f_p^{r\alpha}\norm f_q^{r(1 -
\alpha)}$.
\end{proof}

\begin{method}\label{met:b3:lp:toolkit}
De $L^p$-gereedschapskist, zoals ze hieronder overal gebruikt
wordt: om een identiteit of ongelijkheid voor alle $f \in L^p$ te
bewijzen --- bewijs haar op een \omterm{def:b3:topology:interior}{dichte} klasse ($\mathcal
C_c^\infty$ via \cref{thm:b3:lp:regularization}) en breid uit met
\omterm{def:b3:topology:continuity}{continuïteit} (\cref{thm:b3:complete:extension}, want beide leden
zijn $L^p$-continu); om $f = 0$ te bewijzen, test tegen $\mathcal
C_c^\infty$ (\cref{cor:b3:lp:fundlemma}); om gladheid te winnen,
convolueer; om exponenten te ruilen, Hölder en interpolatie. De
fouriertheorie van \cref{ch:b3:fouriertransform} is één lange
toepassing van deze methode.
\end{method}

\section{Oefeningen}

\begin{exercise}[$\star$]\label{exo:b3:lp:1}
(a) Formuleer en bewijs de ongelijkheid van Cauchy--Schwarz in
$L^2(\mu)$ als het geval $p = q = 2$ van Hölder.
(b) Toon aan dat op een \emph{kansruimte} $p \mapsto \norm f_p$
niet-dalend is.
(c) Wanneer is Hölder een gelijkheid voor $p = 1$, $q = \infty$?
\end{exercise}

\begin{exercise}[$\star$]\label{exo:b3:lp:2}
Voor welke $p \in \intco1\infty$ behoren de volgende functies tot
$L^p$?
\[
x^{-1/2}\mathbf 1_{\intoo01},\qquad
x^{-1/2}\mathbf 1_{\intoo1\infty},\qquad
\frac{1}{x^{1/2}(1 + \abs{\ln x})}\ \text{op } \intoo01,
\qquad
\frac1{1 + \abs x}\ \text{op } \R .
\]
Besluit: op $\intoo01$ is een kleine $p$ makkelijker, op
$\intoo1\infty$ een grote $p$, en op $\R$ bevat geen enkele
$L^p$ een andere.
\end{exercise}

\begin{exercise}[$\star\star$]\label{exo:b3:lp:3}
(De schrijfmachine) Nummer de dyadische intervallen $I_1 =
\intcc01$, $I_2 = \intcc0{\frac12}$, $I_3 = \intcc{\frac12}1$,
$I_4 = \intcc0{\frac14}$, \dots\ en stel $f_n = \mathbf
1_{I_n}$.
(a) Toon aan dat $f_n \to 0$ in elke $L^p(\intcc01)$ met $p <
\infty$, maar dat $(f_n(x))$ voor \emph{elke} $x \in \intcc01$
divergeert.
(b) Geef de b.o.\ convergente deelrij die
\cref{thm:b3:lp:complete} belooft.
(c) Geef omgekeerd een rij die b.o.\ convergeert maar niet in
$L^1$, en één die in $L^1$ convergeert maar in geen enkele $L^p$
met $p > 1$.
\end{exercise}

\begin{exercise}[$\star\star$]\label{exo:b3:lp:4}
Zij $\mu(X) < \infty$ en $f \in L^\infty(\mu)$ met $f \neq 0$.
Toon aan dat $\norm f_p \to \norm f_\infty$ als $p \to \infty$.
\emph{(Bovengrens via (a) van \cref{prop:b3:lp:inclusions};
ondergrens door te integreren over $\{\abs f > \norm f_\infty -
\varepsilon\}$, van positieve \omterm{def:b3:measure:measure}{maat}.)}
\end{exercise}

\begin{exercise}[$\star\star$]\label{exo:b3:lp:5}
(a) Waar precies gebruikt het bewijs van
\cref{thm:b3:lp:density}(3) dat $p < \infty$?
(b) Toon aan dat $f \in L^\infty(\R)$ voldoet aan
$\norm{\tau_hf - f}_\infty \to 0$ dan en slechts dan als $f$ een
uniform \omterm{def:b3:topology:continuity}{continue} vertegenwoordiger heeft.
\end{exercise}

\begin{exercise}[$\star\star$]\label{exo:b3:lp:6}
Zij $p, q$ toegevoegd, $f \in L^p(\R^d)$ en $g \in L^q(\R^d)$.
Toon aan dat $f * g$ \emph{overal} gedefinieerd en begrensd is,
met $\norm{f*g}_\infty \leq \norm f_p\norm g_q$, en bovendien
\emph{uniform \omterm{def:b3:topology:continuity}{continu}}. \emph{(\omterm{def:b3:topology:continuity}{Continuïteit} van de translatie in
$L^p$; behandel $p \in \{1, \infty\}$ apart --- gebruik voor $p =
\infty$ de \omterm{def:b3:topology:continuity}{continuïteit} van de translatie op de
$L^1$-factor.)}
\end{exercise}

\begin{exercise}[$\star\star$]\label{exo:b3:lp:7}
Zij $f \in L^1_{\mathrm{loc}}(\intoo ab)$ met $\int f\varphi' =
0$ voor elke $\varphi \in \mathcal C^\infty_c(\intoo ab)$. Toon
aan dat $f$ b.o.\ gelijk is aan een constante. \emph{(Leg $\chi
\in \mathcal C_c^\infty$ vast met $\int\chi = 1$; elke $\psi \in
\mathcal C_c^\infty$ met $\int\psi = 0$ is een $\varphi'$;
schrijf een algemene testfunctie als $\psi + (\int\psi)\chi$ en
pas \cref{cor:b3:lp:fundlemma} toe op $f - c$ met $c = \int
f\chi$.)}
\end{exercise}

\begin{exercise}[$\star\star\star$]\label{exo:b3:lp:8}
(Gladde Urysohn) Zij $K \subseteq U \subseteq \R^d$ met $K$
\omterm{def:b3:topology:compact}{compact} en $U$ open. Construeer $\varphi \in \mathcal
C^\infty_c(\R^d)$ met $0 \leq \varphi \leq 1$, $\varphi = 1$ op
$K$ en $\operatorname{supp}\varphi \subseteq U$.
\emph{(Strijk de indicator van de $\delta$-omgeving $K_\delta$
van $K$ glad met $\rho_{\delta/2}$, voor kleine $\delta$.)}
Leid er een $\mathcal C^\infty$-partitie van de eenheid uit af
voor een \omterm{def:b3:topology:compact}{compacte} verzameling die door eindig veel \omterm{def:b3:topology:topology}{open
verzamelingen} wordt overdekt.
\end{exercise}

\begin{exercise}[$\star\star$]\label{exo:b3:lp:9}
Met interpolatie (\cref{prop:b3:lp:inclusions}(c)):
(a) toon aan dat $L^1(\R)\cap L^\infty(\R) \subseteq L^p(\R)$
voor alle $p$, met $\norm f_p \leq \norm f_1^{1/p}\norm
f_\infty^{1 - 1/p}$;
(b) toon aan dat $f \mapsto \norm f_p$ voor vaste $f$
log-convex is in $\frac1p$, en geef een voorbeeld waarin $f \in
L^p$ precies voor $p$ in een gegeven interval $(p_0, p_1)$.
\end{exercise}

\begin{exercise}[$\star\star$]\label{exo:b3:lp:10}
(Jensen) Zij $\mu$ een \emph{kansmaat}, $f \in L^1(\mu)$ reëel en
$\Phi \colon \R \to \R$ convex. Toon aan dat
\[
\Phi\Bigl(\int f\,\dd\mu\Bigr) \leq \int \Phi\circ
f\,\dd\mu
\]
\emph{(steunlijn van $\Phi$ in het punt $m = \int f$)}. Leid de
ongelijkheid tussen het rekenkundig en het meetkundig gemiddelde
af, en de monotonie van $p \mapsto \norm f_p$ uit
\cref{exo:b3:lp:1}(b).\index{ongelijkheid van Jensen}
\end{exercise}

\begin{exercise}[$\star\star\star$]\label{exo:b3:lp:11}
(Convolutieongelijkheid van Young) Zij $1 \leq p, q, r \leq
\infty$ met $\frac1p + \frac1q = 1 + \frac1r$, en $f \in
L^p(\R^d)$, $g \in L^q(\R^d)$.
(a) Bewijs $\norm{f * g}_r \leq \norm f_p\,\norm g_q$.
\emph{(Schrijf, met de uit $p, q, r$ berekende toegevoegde
exponenten,
\[
\abs{f(y)g(x-y)} = \bigl(\abs f^p\abs g^q\bigr)^{1/r}
\cdot\abs f^{\,p(1/p - 1/r)}\cdot\abs g^{\,q(1/q - 1/r)},
\]
en pas de ongelijkheid van Hölder met drie factoren toe, met de
exponenten $r$, $\frac{pr}{r - p}$, $\frac{qr}{r - q}$;
integreer daarna naar $x$ met Tonelli.)}
(b) Ga de drie al bekende bijzondere gevallen na: $r = \infty$
(Hölder, \cref{exo:b3:lp:6}); $q = 1$ ($L^p$-stabiliteit van de
convolutie met een \omterm{def:b3:lebesgue:l1}{integreerbare} kern); $p = q = 1$ ($L^1$ is
een convolutiealgebra, \cref{thm:b3:product:convolution}).
(c) Waarom is er geen ongelijkheid met $\frac1p + \frac1q < 1 +
\frac1r$? \emph{(Test op de dilataties $f_\lambda(x) =
f(\lambda x)$ en vergelijk de schalingen van beide leden.)}
\end{exercise}

\begin{exercise}[$\star\star$]\label{exo:b3:lp:12}
(Gelijkheidsgevallen) (a) Toon aan dat in de ongelijkheid van
Hölder $\int\abs{fg} \leq \norm f_p\norm g_q$ ($1 < p < \infty$)
gelijkheid geldt dan en slechts dan als $\abs f^p$ en $\abs g^q$
b.o.\ evenredig zijn. \emph{(Volg het gelijkheidsgeval van de
ongelijkheid van Young $ab \leq \frac{a^p}p + \frac{b^q}q$, dat
$a^p = b^q$ luidt.)}
(b) Toon aan dat in de ongelijkheid van Minkowski $\norm{f + g}_p
\leq \norm f_p + \norm g_p$ ($1 < p < \infty$) gelijkheid met $f,
g \neq 0$ afdwingt dat $g = cf$ b.o., met $c > 0$.
(c) Contrasteer met $p = 1$ en $p = \infty$: beschrijf de (veel
grotere) gelijkheidsgevallen daar, aan de hand van voorbeelden.
\end{exercise}

\section{Probleem: de ongelijkheid van Hardy}

\begin{problem}[{Weekendopgave --- de ongelijkheid van Hardy en
haar scherpe constante}]\label{pb:b3:lp:1}
Voor $f \in L^p(\intoo0{+\infty})$ met $1 < p < \infty$
definiëren we de \emph{operator van Hardy}
\[
(Hf)(x) = \frac1x\int_0^x f(t)\,\dd t .
\]
De ongelijkheid van Hardy (1920) beweert\index{ongelijkheid van
Hardy}
\[
\norm{Hf}_p \;\leq\; \frac{p}{p-1}\,\norm f_p ,
\]
en de constante $\frac p{p-1}$ is optimaal en wordt niet bereikt.
Deze opgave bewijst alles en breidt daarna uit naar reeksen.

\medskip
\noindent\textbf{Deel I --- De ongelijkheid.} Neem eerst $f \geq
0$ \omterm{def:b3:topology:continuity}{continu} met \omterm{def:b3:topology:compact}{compacte} drager in $\intoo0{+\infty}$, en stel
$F(x) = \int_0^xf$.
\begin{enumerate}
  \item Toon aan dat $Hf \in L^p$: nabij $0$ verdwijnt $F$ op een
        \omterm{def:b3:topology:topology}{omgeving} van $0$; nabij $\infty$ is $F$ begrensd, dus
        $(Hf)(x) = O(1/x)$, en $x \mapsto \frac1x$ behoort tot
        $L^p(\intoo1{+\infty})$ voor $p > 1$.
  \item Integreer partieel om aan te tonen dat
        \[
        \int_0^\infty \Bigl(\frac Fx\Bigr)^p\dd x
        = \frac{p}{p-1}\int_0^\infty\Bigl(\frac
        Fx\Bigr)^{p-1}f(x)\,\dd x .
        \]
        \emph{(Differentieer $x^{1-p}F^p$; de randtermen
        verdwijnen --- rechtvaardig beide uiteinden.)}
  \item Pas Hölder toe op het rechterlid en leid
        $\norm{Hf}_p \leq \frac p{p-1}\norm f_p$ af voor
        dergelijke $f$.
  \item Breid uit tot heel $L^p$: construeer voor $f \geq 0$
        functies $f_n$, \omterm{def:b3:topology:continuity}{continu} met \omterm{def:b3:topology:compact}{compacte} drager in
        $\intoo0{+\infty}$, met $0 \leq f_n \nearrow f$ b.o.\
        \emph{(knot af en benader daarna monotoon ---
        rechtvaardig de constructie)}; dan is $Hf_n \nearrow Hf$
        puntsgewijs (monotone convergentie binnen het
        gemiddelde) en draagt de monotone convergentie de
        ongelijkheid over op de limiet. Besluit voor getekende
        of complexe $f$ met $\abs{Hf} \leq H\abs f$.
\end{enumerate}

\medskip
\noindent\textbf{Deel II --- Optimaliteit.}
\begin{enumerate}[resume]
  \item Stel voor $A > 1$ $f_A(t) = t^{-1/p}\,\mathbf
        1_{\intcc1A}(t)$. Bereken $\norm{f_A}_p^p = \ln A$ en,
        voor $1 \leq x \leq A$,
        \[
        (Hf_A)(x) = \frac p{p-1}\;x^{-1/p}\,
        \bigl(1 - x^{-(1 - 1/p)}\bigr).
        \]
  \item Leid $\liminf_{A\to\infty} \norm{Hf_A}_p/
        \norm{f_A}_p \geq \frac{p}{p-1}$ af en besluit dat de
        constante optimaal is.
  \item Toon aan dat gelijkheid $\norm{Hf}_p = \frac
        p{p-1}\norm f_p$ met $f \neq 0$ onmogelijk is.
        \emph{(Volg het gelijkheidsgeval van Hölder in vraag 3:
        het zou een gedrag $f = cx^{-1/p}$ afdwingen, dat niet
        in $L^p$ ligt.)}
\end{enumerate}

\medskip
\noindent\textbf{Deel III --- De discrete ongelijkheid.}
\begin{enumerate}[resume]
  \item Vergelijk voor een niet-stijgende $g \geq 0$ op
        $(0,\infty)$ en $a_n = g(n)$ de grootheden $\sum a_n^p$
        en $\int g^p$, en de bijbehorende $H$-gemiddelden, om
        uit Deel I de \emph{discrete ongelijkheid van Hardy}
        af te leiden: voor $a_n \geq 0$ is
        \[
        \sum_{n\geq1}\Bigl(\frac{a_1 + \dots +
        a_n}{n}\Bigr)^{p}
        \;\leq\;
        \Bigl(\frac{p}{p-1}\Bigr)^{p}\,\sum_{n\geq1}a_n^p
        \]
        --- bewijs haar eerst voor niet-stijgende $(a_n)$ via de
        bovenstaande vergelijking, en herleid daarna het
        algemene geval tot het niet-stijgende met een
        herschikking \emph{(neem, met een rechtvaardiging van
        één regel, aan dat $(a_n)$ dalend sorteren het
        linkerlid alleen kan vergroten terwijl het rechterlid
        onveranderd blijft)}.
  \item Leid af: is $\sum a_n^p < \infty$, dan liggen de
        cesàrogemiddelden van $(a_n)$ opnieuw in $\ell^p$ --- en
        geef een voorbeeld ($p = 2$) waarin $(a_n) \in \ell^2$
        maar $a_n$ niet sommeerbaar is, terwijl Hardy de
        gemiddelden toch beheerst.
\end{enumerate}

\medskip
\noindent\textbf{Deel IV --- Epiloog.}
\begin{enumerate}[resume]
  \item Toon aan dat de ongelijkheid van Hardy faalt voor $p =
        1$: bereken $Hf$ voor $f = \mathbf 1_{\intcc01}$ en
        merk op dat $Hf \notin L^1$. Waar breekt het bewijs?
\end{enumerate}

\medskip
\noindent\textbf{Deel V --- De maximaalfunctie en de
differentiatiestelling van Lebesgue.} Hardy middelt vanaf de
oorsprong; Hardy en Littlewood middelen \emph{rond elk punt}.
Voor $f \in L^1(\R)$ definiëren we
\[
Mf(x) = \sup_{r>0}\ \frac1{2r}\int_{x-r}^{x+r}\abs
f\,\dd\lambda .
\]
\begin{enumerate}[resume]
  \item (Vitali, eindige versie) Zij $B_1, \dots, B_N$ open
        intervallen. Toon aan dat er een disjuncte deelfamilie
        $B_{i_1}, \dots, B_{i_k}$ bestaat met $\bigcup_jB_j
        \subseteq \bigcup_l3B_{i_l}$, waarbij $3B$ het interval
        is met hetzelfde middelpunt en de drievoudige lengte
        \emph{(gulzig: kies telkens het langste interval dat
        disjunct is van de reeds gekozen)}.
  \item (Zwak type $(1,1)$) Toon aan dat voor elke $t > 0$
        \[
        \lambda\bigl(\{Mf > t\}\bigr) \;\leq\;
        \frac3t\,\norm f_1 :
        \]
        elke $x$ met $Mf(x) > t$ bezit een gecentreerd interval
        $B_x$ met $\int_{B_x}\abs f > t\,\lambda(B_x)$; neem een
        \omterm{def:b3:topology:compact}{compacte} $K \subseteq \{Mf > t\}$ (inwendige
        regulariteit), overdek haar met eindig veel $B_x$, pas
        vraag 11 toe en put uit.
  \item Bereken $M\mathbf 1_{\intcc01}(x)$ voor $x > 1$ en leid
        af dat $Mf \notin L^1$ voor elke $f \neq 0$ ($Mf(x) \geq
        \frac c{\abs x}$ op oneindig): bij $p = 1$ is de zwakke
        ongelijkheid van vraag 12 de best mogelijke uitspraak.
  \item (Sterk type voor $p > 1$) Voor $f \in L^p$: splits $f =
        f\,\mathbf 1_{\abs f > t/2} + f\,\mathbf 1_{\abs f
        \leq t/2}$, merk op dat $Mf \leq M\bigl(f\mathbf 1_{\abs
        f > t/2}\bigr) + \frac t2$, en combineer vraag 12 met de
        laagjesformule (\cref{prop:b3:product:layercake}) en
        Tonelli om
        \[
        \norm{Mf}_p^p \;\leq\;
        \frac{6p\,2^{p-1}}{p-1}\,\norm f_p^p
        \]
        te bewijzen. (De explosie als $p \downarrow 1$ is het
        falen van vraag 13, gekwantificeerd.)
  \item (Differentiatiestelling van Lebesgue) Bewijs: voor $f
        \in L^1(\R)$ is
        \[
        \frac1{2r}\int_{x-r}^{x+r}f\,\dd\lambda
        \;\longrightarrow\; f(x)
        \qquad (r \to 0)\quad\text{voor b.o.\ }x .
        \]
        \emph{(Duidelijk voor \omterm{def:b3:topology:continuity}{continue} $f$. Schrijf in het
        algemeen $f = g + h$ met $g$ \omterm{def:b3:topology:continuity}{continu} met \omterm{def:b3:topology:compact}{compacte} drager
        en $\norm h_1 < \varepsilon$
        (\cref{thm:b3:lp:density}); de verzameling waar de
        $\limsup_{r\to0}$ van de gemiddelde oscillatie $\delta$
        overtreft, ligt binnen $\{Mh > \delta/2\} \cup \{\abs h
        > \delta/2\}$, van \omterm{def:b3:measure:measure}{maat} $O(\varepsilon/\delta)$; laat
        $\varepsilon \to 0$ gaan, en daarna $\delta \to 0$ langs
        een rij.)}
  \item Leid af: (a) bijna elk punt is een \emph{lebesguepunt}
        van $f$; (b) voor $f \in L^1$ is de primitieve $F(x) =
        \int_0^xf$ b.o.\ differentieerbaar met $F' = f$ b.o.\
        --- de integraalhelft van de hoofdstelling van de
        integraalrekening in de wereld van Lebesgue, die de
        cirkel sluit die de trapfunctie
        (\cref{pb:b3:measure:1}) opende door te tonen dat de
        andere helft kan falen.
  \item (Dichtheidspunten) Toon voor \omterm{def:b3:lebesgue:measurable}{meetbare} $A \subseteq \R$
        aan dat bijna elke $x \in A$ voldoet aan
        $\frac{\lambda(A\cap\intcc{x-r}{x+r})}{2r} \to 1$:
        \omterm{def:b3:lebesgue:measurable}{meetbare} verzamelingen zijn lokaal vol in bijna al hun
        punten. Schets in twee regels hoe dit alweer een ander
        bewijs van de stelling van Steinhaus
        (\cref{exo:b3:measure:8}) oplevert.
\end{enumerate}

\medskip
\noindent\textbf{Deel VI --- Variaties op het thema van het
middelen.}
\begin{enumerate}[resume]
  \item (Gewogen Hardy) Toon voor $\alpha < p - 1$ en $f \geq 0$
        aan dat
        \[
        \int_0^\infty\Bigl(\frac{F(x)}x\Bigr)^{p}
        x^{\alpha}\,\dd x \;\leq\;
        \Bigl(\frac{p}{p - 1 - \alpha}\Bigr)^{p}
        \int_0^\infty f(x)^p\,x^{\alpha}\,\dd x
        \]
        met dezelfde partiële integratie, en ga na dat het
        grensgeval $\alpha = p - 1$ werkelijk verboden is (pas
        het tegenvoorbeeld van vraag 10 aan).
  \item (De toegevoegde operator) Zij $H^*f(x) =
        \int_x^{\infty}\frac{f(t)}t\,\dd t$. Toon aan dat
        $\langle Hf, g\rangle = \langle f, H^*g\rangle$ voor
        niet-negatieve $f, g$ (Tonelli), en bewijs
        $\norm{H^*f}_p \leq p\,\norm f_p$ \emph{(rechtstreeks
        met partiële integratie, of uit Hardy op de toegevoegde
        exponent via dualiteit --- let op welke exponent welke
        constante oplevert)}.
  \item (Een ongelijkheid van hilberttype) Leid af dat voor
        niet-negatieve $f \in L^p$, $g \in L^q$
        \[
        \int_0^\infty\!\!\int_0^\infty
        \frac{f(x)\,g(y)}{\max(x,y)}\,\dd x\,\dd y
        \;\leq\; (p + q)\,\norm f_p\,\norm g_q
        \]
        \emph{(splits langs $y < x$ / $y \geq x$: elke helft is
        een paring van de ene functie tegen een
        Hardy-getransformeerde van de andere)}.
  \item (Optimaliteit, discreet) Toon aan dat ook de constante
        $\bigl(\frac p{p-1}\bigr)^p$ van vraag 8 optimaal is:
        test op $a_n = n^{-1/p}\,\mathbf 1_{n \leq N}$,
        vergelijk beide leden met integralen en laat $N \to
        \infty$ gaan (de discrete spiegel van Deel II).
  \item (Synthese) Drie middelingsoperatoren zijn in deze opgave
        opgedoken: die van Hardy, $H$, het discrete
        cesàrogemiddelde en de maximaaloperator $M$. Formuleer
        in één regel per stuk wat hun begrensdheid zegt, merk op
        dat alle drie precies bij $p = 1$ falen, en verklaar
        waarom het driemaal hetzelfde falen is (de harmonische
        staart $\frac1x$).
\end{enumerate}

\medskip
\noindent\textbf{Deel VII --- De ongelijkheid van Carleman, en
hoe scherp scherp is.}
\begin{enumerate}[resume]
  \item (Ongelijkheid van Carleman) Zij $a_n \geq 0$ met $\sum
        a_n < \infty$. Pas de discrete ongelijkheid van Hardy
        van vraag 8 toe op $b_n = a_n^{1/p}$, gebruik de
        ongelijkheid tussen het rekenkundig en het meetkundig
        gemiddelde en laat $p \to \infty$ gaan (toon aan dat $p
        \mapsto \bigl(\frac p{p-1}\bigr)^p$ daalt naar $\eu$) om
        \[
        \sum_{n\geq1}\bigl(a_1a_2\cdots a_n\bigr)^{1/n}
        \;\leq\; \eu\,\sum_{n\geq1}a_n
        \]
        te verkrijgen: de meetkundige gemiddelden van een
        sommeerbare rij zijn opnieuw sommeerbaar, tegen een prijs
        van hoogstens $\eu$.
  \item (De constante $\eu$ is optimaal) Test $a_n =
        \frac1n\,\mathbf 1_{n\leq N}$: toon met de
        Stirling-insluiting van \cref{pb:b3:product:1} aan dat
        $(n!)^{-1/n} = \frac\eu n\bigl(1 + O\bigl(\frac{\ln
        n}n\bigr)\bigr)$, leid af dat beide leden van Carleman
        als $\eu\ln N$ groeien, en besluit dat geen enkele
        constante kleiner dan $\eu$ kan werken. (Merk het
        patroon op: de optimalisatoren van Hardy en van Carleman
        zijn beide de rijen van harmonisch type die er net niet
        in slagen in de ruimte te liggen.)
  \item (Hoe traag wordt ``scherp'' benaderd?) Neem $p = 2$.
        Bereken voor $f = \mathbf 1_{\intcc01}$ de verhouding
        $\norm{Hf}_2/\norm f_2 = \sqrt2$, tegenover de grens
        $2$. Bewijs voor de bijna-optimalisatoren $f_A$ van
        vraag 5 de exacte identiteit
        \[
        \norm{Hf_A}_2^2 = 4\ln A - 8 + \frac{8}{\sqrt A},
        \qquad\text{dus}\qquad
        \frac{\norm{Hf_A}_2^2}{\norm{f_A}_2^2}
        = 4 - \frac{8 - 8A^{-1/2}}{\ln A} .
        \]
        Evalueer bij $A = \eu^{10}$ (verhouding $\approx 1.790$)
        en becommentarieer: het supremum $2$ wordt slechts met
        snelheid $1/\ln A$ benaderd --- een optimale constante
        kan numeriek zo goed als onzichtbaar zijn.
\end{enumerate}
\end{problem}
